
\subsubsection{5.1 M1: Cluster Formation}

SimCLR training for 20 epochs on Waterbirds produced embeddings with the clusterability metrics shown in Table 1.

\textbf{Table 1: Embedding Clusterability Metrics for SimCLR (20 epochs, ResNet-50)}

\begin{table}[htbp]
\centering
\resizebox{\columnwidth}{!}{%
\begin{tabular}{llll}
\toprule
Metric & Value & Threshold & Result \\
\midrule
AMI & 0.2795 & $\geq$ 0.4 & Below threshold \\
Silhouette Score & 0.2967 & $\geq$ 0.3 & Below threshold \\
Spurious Correlation & 93\% & - & - \\
\bottomrule
\end{tabular}
}
\end{table}

AMI was 0.2795, below the 0.4 threshold. Silhouette score was 0.2967, marginally below the 0.3 threshold. Under the pre-registered success criteria, M1 was not supported at 20 epochs of training.

AMI values at epoch 10 and epoch 20 were 0.2786 and 0.2795 respectively, showing minimal change over this interval.

\subsubsection{5.2 M2: Clusterability-Efficacy Relationship}

To test whether AMI predicts cluster-balanced retraining efficacy, models were stratified by AMI threshold and $\Delta WGA$ was computed. With n=2 samples (epoch 10 and epoch 20 checkpoints), statistical tests have limited power.

\textbf{Table 2: AMI-Efficacy Relationship}

\begin{table}[htbp]
\centering
\resizebox{\columnwidth}{!}{%
\begin{tabular}{llll}
\toprule
Stratum & Mean AMI & Mean $\Delta WGA$ (pp) & Sample Count \\
\midrule
AMI $\geq$ 0.28 & 0.285 & 0.00 & 1 \\
AMI < 0.28 & 0.279 & -5.14 & 1 \\
\bottomrule
\end{tabular}
}
\end{table}

Pearson correlation between AMI and $\Delta WGA$ was r = -1.0 with p = 1.000. Neither stratum achieved the hypothesized $\geq 2.0pp$ improvement for high-AMI models. The small sample size (n=2) limits the interpretability of correlation statistics; a minimum of approximately 20 samples is typically required for reliable correlation estimates.

Under the pre-registered success criteria (positive correlation with p < 0.05, high-AMI mean $\Delta WGA \geq$ 2pp), M2 was not supported at 20 epochs of training.

\subsubsection{5.3 M3: LA-SSL Geometric Effects}

Table 3 compares clusterability and linear separability between SimCLR and LA-SSL at 20 epochs.

\textbf{Table 3: Geometric Comparison Between SimCLR and LA-SSL (20 epochs)}

\begin{table*}[htbp]
\centering
\resizebox{\dimexpr 2\columnwidth+\columnsep\relax}{!}{%
\begin{tabular}{llllll}
\toprule
Metric & SimCLR & LA-SSL & Change & Threshold & Result \\
\midrule
AMI & 0.2795 & 0.2852 & +2.04\% & $\geq 30$\% reduction & Opposite direction \\
Linear AUC & 0.9802 & 0.9856 & +0.54\% & $\Delta$ < 5\% & Within threshold \\
\bottomrule
\end{tabular}
}
\end{table*}

LA-SSL produced AMI = 0.2852, representing a 2.04\% relative increase compared to SimCLR's AMI = 0.2795 (absolute increase of 0.0057). The hypothesized $\geq 30$\% reduction in AMI was not observed; instead, AMI increased. Linear probe AUC remained high for both methods (SimCLR: 0.9802, LA-SSL: 0.9856), with $\Delta AUC$ = 0.0054, indicating that linear separability was preserved.

Under the pre-registered success criteria (AMI reduction $\geq 30$\% and $\Delta AUC$ <0.05), M3 was not supported. While linear separability was preserved, the clusterability change was in the opposite direction from the hypothesis.

\subsubsection{5.4 Linear Separability Despite Low Clusterability}

Linear probe AUC for subgroup classification was approximately 0.98 for both SimCLR and LA-SSL despite AMI values of approximately 0.28. This indicates that linear separability and discrete clusterability are dissociable properties under the tested conditions.

\subsubsection{5.5 Implementation Validation}

To verify that results were not due to implementation errors, comprehensive code validation was performed. The h-e1 implementation passed 43/43 unit and integration tests (100\% pass rate) with 100\% software development document compliance across 15 tasks and 0 critical or blocking issues. The h-m-integrated implementation passed 5/5 tests with 100\% SDD compliance across 43 tasks and 0 critical or blocking issues.
